\subsection{Contraction of the arrows of a homotopy}

Let H be a quiver, let G be a groupoid, let \(\varphi\) and \(\psi\) be morphisms of quivers from H to G, and let \(h\colon \mathrm{Som}(\mathrm{H})\to \mathrm{Fl}(\mathrm{G})\) be a homotopy connecting \(\varphi\) to \(\psi\). Let G' be the groupoid obtained from G by contracting the arrows in the image of \(h\) (II, p. 175, no. 11), and let \(\beta\colon\mathrm{G}\to\mathrm{G}'\) be the canonical morphism.

Denote by \(\Gamma\) the quiver (Som(G), Som(H), Som(φ), Som(ψ)). By definition of a homotopy, the pair (Id \(_{\text{Som(G)}}\) , \(h\) ) is a quiver morphism from \(\Gamma\) to G; it extends uniquely to a groupoid morphism \(\eta\) : Grp(Γ) → G. By construction, the set of vertices of G' is the set of connected components of the quiver \(\Gamma\) .

Throughout the rest of this section, \(n^{o}\), we shall assume that the groupoid G is transitive. By Remark 1 of II, p. 170, this is equivalent to assuming that the groupoid G' is transitive. We also fix a vertex \(a_{0}\) of G.

Denote by \(\widetilde{\Gamma}\) the graph associated with the quiver \(\Gamma\) (cf. II, p. 156). The set of loops of length \(\geqslant 1\) in \(\widetilde{\Gamma}\) is identified with the set \(\Omega(\widetilde{\Gamma})\) of finite sequences \((z_1, \ldots, z_n)\) , where \(n\) is an integer such that \(n \geqslant 1\) and \(z_1, \ldots, z_n\) are elements of \(\mathrm{Fl}(\widetilde{\Gamma})\) such that \(t(z_i) = o(z_{i+1})\) for every integer \(i\) such that \(1 \leqslant i < n\) and \(t(z_n) = o(z_1)\) .

Let \(\mathbf{z} = ((b_{1}, \varepsilon_{1}), \ldots, (b_{n}, \varepsilon_{n}))\) be an element of \(\Omega(\widetilde{\Gamma})\) . The conjugacy class of the element

\[
\operatorname{Int} (g) (\eta (\mathbf {z})) = g h (b _ {1}) ^ {\varepsilon_ {1}} \dots h (b _ {n}) ^ {\varepsilon_ {n}} g ^ {- 1}
\]

in \(G_{a_{0}}\) does not depend on the choice of the arrow g of G connecting the vertex \(a_{0}\) to the origin of \((b_{1},\varepsilon)\) . We denote by \(c(\mathbf{z})\) this conjugacy class.

PROPOSITION 1. --- The group homomorphism \(\beta_{a_0} \colon \mathrm{G}_{a_0} \to \mathrm{G}'_{\beta(a_0)}\) is surjective, and its kernel is the smallest subgroup of \(\mathrm{G}_{a_0}\) containing the conjugacy classes \(c(\mathbf{z})\) for \(\mathbf{z} \in \Omega(\widetilde{\Gamma})\) .

Let K be the smallest distinguished subgroupoid of G whose set of arrows contains the image of h. The morphism \(\mathrm{G}_{a_{0}} \rightarrow \mathrm{G}^{\prime}_{\beta(a_{0})}\) is surjective and its kernel is equal to \(K_{a_{0}}\) (II, p. 170, prop. 2). The proposition then follows from prop. 8 of II, p. 176.

DEFINITION 1. --- One says that a subset Z of \(\Omega(\widetilde{\Gamma})\) is distinguished (relative to the pair \((\varphi, \psi)\) ) if it satisfies the following properties:

\begin{enumerate}
\def\labelenumi{(\roman{enumi})}
\item
  For every arrow z of Γ, one has \((z,\overline{z})\in\mathbb{Z}\) ;
\item
  For every \((z_{1},\ldots ,z_{n})\in \mathbb{Z}\) , we have \((\overline{z}_n,\dots ,\overline{z}_2,\overline{z}_1)\in \mathbb{Z}\) and \((z_{n},z_{1},\ldots ,z_{n - 1})\in \mathbb{Z}\) ;
\item
  Let \(\mathbf{z} = (z_1, \ldots, z_n)\) and \(\mathbf{z}' = (z_1', \ldots, z_m')\) be elements of \(Z\) such that \(t(z_n) = o(z_1')\) . Set \(\mathbf{z}\mathbf{z}' = (z_1, \ldots, z_n, z_1', \ldots, z_m')\) . If two elements among z, z', zz' belong to Z, then so does the third;
\item
  For every arrow \(f\) of H, put \(\widetilde{\varphi}(f,1) = \widetilde{\psi}(f,-1) = \varphi(f)\) and \(\widetilde{\varphi}(f,-1) = \widetilde{\psi}(f,1) = \psi(f)\) . Let \(n\) be an integer \(\geqslant 1\) and \((f_1,\varepsilon_1),\ldots,(f_n,\varepsilon_n)\) a sequence of elements of \(\mathrm{Fl}(\mathrm{H})\times \{-1,1\}\) such that \(\widetilde{\psi}(f_i,\varepsilon_i) = \widetilde{\varphi}(f_{i+1},\varepsilon_{i+1})\) for \(1\leqslant i < n\) and \(\widetilde{\psi}(f_n,\varepsilon_n) = \widetilde{\varphi}(f_1,\varepsilon_1)\) ; let us denote by \(a_i\) the origin of \(f_i\) and \(b_i\) its term. For \(((a_1,\varepsilon_1),\ldots,(a_n,\varepsilon_n))\) to belong to Z, it is necessary and sufficient that \(((b_1,\varepsilon_1),\ldots,(b_n,\varepsilon_n))\) belong to Z.
\end{enumerate}

The intersection in \(\Omega(\widetilde{\Gamma})\) of every family of distinguished subsets (relative to \((\varphi, \psi)\) ) is again distinguished. In particular, there exists a smallest distinguished subset containing a given subset \(Z\) of \(\Omega(\widetilde{\Gamma})\) .

PROPOSITION 2. --- If \(N\) is a normal subgroup of \(G_{a}\), the set of elements \(\mathbf{z} \in \Omega(\widetilde{\Gamma})\) such that \(c(\mathbf{z})\) is contained in \(N\) is a distinguished subset of \(\Omega(\widetilde{\Gamma})\).

Let Z denote the set of elements \(\mathbf{z} \in \Omega(\widetilde{\Gamma})\) such that \(c(\mathbf{z}) \subset \mathrm{N}\) .

Let \(z \in \mathrm{Fl}(\widetilde{\Gamma})\) . We have \(c(z, \overline{z}) = \{e_a\}\) by definition, whence \((z, \overline{z}) \in \mathbf{Z}\) . Let \((z_1, \ldots, z_n) \in \mathbf{Z}\) . By definition of \(c\) , the conjugacy class \(c(z_1, \ldots, z_n)\) is equal to \(c(z_n, z_1, \ldots, z_{n-1})\) and is formed by the inverses of the elements of \(c(\overline{z}_n, \ldots, \overline{z}_1)\) . This shows that \(\mathbf{Z}\) satisfies condition (ii).

With the notations of condition (iii), one can choose elements \(u \in c(\mathbf{z})\) , \(v \in c(\mathbf{z}')\) and \(w \in c(\mathbf{z}\mathbf{z}')\) such that uv = w. If two of the elements u, v, and w belong to N, then so does the third, hence N satisfies (iii).

Using the notation from (iv), for every integer i such that \(1 \leqslant i \leqslant n\) , the relation

\[
  \varphi (f _ {i}) h (b _ {i}) = h (a _ {i}) \psi (f _ {i}),
\]

since h is a homotopy connecting \(\varphi\) to \(\psi\). This equality can also be written as

\[
\widetilde {\varphi} (f _ {i}, \varepsilon_ {i}) h (b _ {i}) ^ {\varepsilon_ {i}} = h (a _ {i}) ^ {\varepsilon_ {i}} \widetilde {\psi} (f _ {i}, \varepsilon_ {i}).
\]

Taking into account the relations \(\widetilde{\psi}(f_i,\varepsilon_i) = \widetilde{\varphi}(f_{i + 1},\varepsilon_{i + 1})\) for \(1\leqslant i <   n\) and \(\widetilde{\psi} (f_n,\varepsilon_n) = \widetilde{\varphi} (f_1,\varepsilon_1)\) , we deduce

\[
\widetilde {\varphi} (f _ {1}, \varepsilon_ {1}) h (b _ {1}) ^ {\varepsilon_ {1}} \dots h (b _ {n}) ^ {\varepsilon_ {n}} = h (a _ {1}) ^ {\varepsilon_ {1}} \dots h (a _ {n}) ^ {\varepsilon_ {n}} \widetilde {\varphi} (f _ {1}, \varepsilon_ {1}),
\]

so that the conjugacy classes \(c((a_{1},\varepsilon_{1}),\ldots,(a_{n},\varepsilon_{n}))\) and \(c((b_{1},\varepsilon_{1}),\ldots,(b_{n},\varepsilon_{n}))\) are equal. This shows that Z satisfies condition (iv) and completes the proof of the proposition.

COROLLARY. --- Let Z be a subset of \(\Omega(\widetilde{\Gamma})\). For the conjugacy classes \(c(\mathbf{z})\), for \(\mathbf{z} \in \mathbf{Z}\), to generate the kernel of the canonical homomorphism \(\beta_{a_0}: \mathrm{G}_{a_0} \to \mathrm{G}'_{\beta(a_0)}\), it suffices that the smallest distinguished subset of \(\Omega(\widetilde{\Gamma})\) containing Z be equal to \(\Omega(\widetilde{\Gamma})\).

Let \(N\) be the smallest subgroup of \(G_{a_{0}}\) containing \(c(\mathbf{z})\) for all \(z \in Z\); it is normal. Let \(Z'\) be the set of elements \(z\) of \(\Omega(\widetilde{\Gamma})\) such that \(c(\mathbf{z}) \in \mathrm{N}\). By Proposition 2, \(Z'\) is a distinguished subset of \(\Omega(\widetilde{\Gamma})\). It contains Z. By hypothesis, it is therefore equal to \(\Omega(\widetilde{\Gamma})\). It then follows from prop. 1 (II, p. 197) that N is the kernel of the homomorphism \(\beta_{a}\).

